%% source: 00_abstract (402w) -> target 200w

\begin{abstract}
Vector tiles are the dominant delivery format for interactive web cartography, and the client redraws each visible tile inside a per-frame budget that covers network transfer, decoding, per-feature style evaluation, and draw submission.
Cartographic styling and tile production are normally handled by separate tools, so the coupling between what a style actually reads and what a tile actually carries stays unexploited.
This article treats styling and tile encoding as a single optimization problem for MapLibre rendering pipelines.
We build an optimizer that rewrites style expressions to a fixpoint, restructures the typed style document, and then derives from the optimized style a pruning advisory that removes unused attributes, geometry types, and feature values before re-encoding each tile in the columnar \ac{MLT} format.
Measured over a corpus of production and institutional styles under interactive navigation scenarios, discarding tile content the style never accounts for most of the benefit, lowering median load time by 71\% and raising sustained frame rate to a median of \num{2.5}$\times$.
The columnar encoder alone shrinks tile volume by 28\% against the uncompressed reference and 12\% against gzip-compressed Mapbox Vector Tiles.
Style rewriting shrinks unseen styles by a median of 30.8\% with visually equivalent output.
Quantifying how the style-level and data-level passes interact, we find their combined effect is additive rather than synergistic.
Per-scenario synergy is present but too weak to raise the aggregate.
\end{abstract}

\begin{keywords}
vector tiles; MapLibre; map style optimization; columnar tile encoding
\end{keywords}